\documentclass[11pt]{article}

\usepackage[utf8]{inputenc}
\usepackage[T1]{fontenc}
\usepackage{amsmath,amssymb}
\usepackage{graphicx}
\usepackage{booktabs}
\usepackage{hyperref}
\usepackage{xcolor}
\usepackage[margin=2.5cm]{geometry}
\usepackage{caption}
\usepackage{subcaption}
\usepackage{float}
\usepackage{enumitem}
\usepackage{multirow}

\hypersetup{colorlinks=true, linkcolor=blue!60!black, citecolor=green!50!black, urlcolor=blue!70!black}

\title{\textbf{Reproducing 1.58-Bit FLUX:\\Ternary Quantization with LoRA Compensation\\for Diffusion Transformers}}

\author{Ugon For \\ \texttt{github.com/ugonfor/1.58bit-flux}}
\date{March 2026}

\begin{document}
\maketitle

\begin{abstract}
We reproduce the core approach of ``1.58-bit FLUX'' (Yang et al., 2024), demonstrating that ternary ($\{-1, 0, +1\}$) quantization of the FLUX.1-dev diffusion transformer can retain up to 90.4\% of BF16 image quality when combined with rank-128 LoRA compensation and offline flow-matching distillation.
Starting from naive post-training quantization that produces pure noise (CLIP 0.178 vs BF16's 0.322), we develop a systematic pipeline recovering image quality through progressive data scaling and architectural improvements over 13 experimental iterations.
Key findings: (1)~a log$_2$ data scaling law accurately predicts quality gains up to 2{,}132 training prompts; (2)~this law breaks at $\sim$4{,}000 prompts due to rank-64 LoRA capacity saturation at $\sim$90\%; (3)~rank-128 LoRA breaks through this ceiling to 90.4\%; (4)~with current INT8 weight storage, inference peak VRAM is reduced by 30.8\% (33.83$\to$23.41\,GB for rank-64) at a 23\% speed overhead; (5)~true 2-bit packed storage requires fused custom kernels not yet available for BF16 activations on current hardware.
All experiments run on a single NVIDIA A100-SXM4-80GB GPU.
\end{abstract}

\section{Introduction}

The FLUX.1-dev transformer \cite{flux} contains 11.9B parameters across 504 linear layers, consuming 22.7\,GB in BF16.
Yang et al.\ \cite{yang2024} propose extreme quantization to 1.58 bits---ternary $\{-1, 0, +1\}$ weights with per-channel scaling and LoRA adapters---claiming near-lossless quality at 7{,}232 training prompts.
Their code and weights have not been released.

This report documents our complete reproduction from scratch: 13 experiments over 10 weeks, from catastrophic PTQ failure to 90.4\% BF16 quality retention.
We provide all code, checkpoints, and training data at \url{https://huggingface.co/ugonfor/1.58bit-flux-reproduction}.

\paragraph{Contributions.}
\begin{itemize}[nosep,leftmargin=1.5em]
  \item Complete open-source ternary FLUX reproduction with flow-matching distillation
  \item Discovery of a log$_2$ data scaling law and its capacity-limited breakdown
  \item Rank-128 LoRA breaking the rank-64 ceiling (90.0\%$\to$90.4\%)
  \item Comprehensive memory/speed benchmarks including attempted custom kernel optimization
  \item Honest assessment: without fused ternary matmul kernels, memory gains are limited to INT8 storage (30.8\% reduction), not the theoretical 8$\times$ from 2-bit packing
\end{itemize}

\section{Method}

\subsection{Ternary Quantization with LoRA}

Each BF16 weight $w$ is quantized via absmean thresholding:
\begin{equation}
  \hat{w}_i = \text{sign}(w_i) \cdot \mathbf{1}[|w_i| > \text{mean}(|w|)]
\end{equation}
The effective weight combines ternary, per-channel scale $s$, and low-rank adapter:
\begin{equation}
  W_{\text{eff}} = \hat{W} \cdot \text{diag}(s) + AB
\end{equation}
where $A \in \mathbb{R}^{d_{\text{out}} \times r}$, $B \in \mathbb{R}^{r \times d_{\text{in}}}$ are initialized via truncated SVD of the quantization residual.
The forward pass avoids materializing the full effective weight:
\begin{equation}
  y = (x \hat{W}^\top) \odot s + (xB^\top)A^\top + b
\end{equation}

\subsection{Offline Flow-Matching Distillation}

We minimize the velocity-field divergence between student and teacher:
\begin{equation}
  \mathcal{L} = \mathbb{E}_{z_0, \epsilon, t} \left[ \| v_\text{s}(z_t, t, c) - v_\text{t}(z_t, t, c) \|^2 \right] + \lambda \cdot \mathcal{L}_{\text{LPIPS}}
\end{equation}
where $z_t = (1-t)z_0 + t\epsilon$, teacher latents $z_0$ are pre-generated via 28-step BF16 inference, $t \sim \text{LogitNormal}(\sigma=0.5)$, and $\lambda = 0.1$.

\subsection{Training Pipeline}

\begin{enumerate}[nosep,leftmargin=1.5em]
  \item \textbf{Teacher dataset}: Generate BF16 images at 1024$\times$1024 (28 steps), store latent $z_0$ + prompt
  \item \textbf{Quantize}: Replace all 504 nn.Linear with TernaryLinear (absmean + SVD-init LoRA)
  \item \textbf{Distill}: Train scale + LoRA via FM loss + LPIPS; AdamW, cosine LR, gradient checkpointing
\end{enumerate}

\section{Experiments}

\subsection{Progression Overview}

Table~\ref{tab:progression} summarizes all 13 experimental iterations.

\begin{table}[h]
\centering
\caption{Experimental progression. OOD CLIP is averaged over 20 unseen diverse prompts as \% of BF16.
$^*$Evaluated on training prompts only (leaky benchmark).}
\label{tab:progression}
\small
\begin{tabular}{llcccc}
\toprule
\textbf{Version} & \textbf{Key Change} & \textbf{Prompts} & \textbf{Rank} & \textbf{Steps} & \textbf{OOD CLIP} \\
\midrule
PTQ & No training & --- & --- & 0 & 55.2\% \\
V1 & FM distillation (cold) & 50 & 64 & 3k & 82.4\%$^*$ \\
V2 & Warm-start + 200 imgs & 200 & 64 & 3k & 95.6\%$^*$ \\
V5 & 174 unique prompts & 174 & 64 & 3k & 101.9\%$^*$ \\
\midrule
V6 & +LPIPS, OOD eval & 174 & 64 & 2k & 86.0\% \\
V7 & 1k prompts, balanced & 1{,}002 & 64 & 4k & 88.9\% \\
V9b & 2k prompts & 2{,}132 & 64 & 6k & \textbf{90.0\%} \\
V9c & 4k prompts (ceiling) & 4{,}007 & 64 & 8k & 88.8\% \\
\midrule
V10 & Rank-128 (cold) & 2{,}132 & 128 & 6k & 87.7\% \\
V10b & Rank-128 (extended) & 2{,}132 & 128 & 12k & \textbf{90.4\%} \\
\bottomrule
\end{tabular}
\end{table}

\subsection{Failed Approaches}

\textbf{Naive PTQ}: Produces pure noise (CLIP 0.178).
\textbf{Layer-wise MSE}: No gradient through output projection.
\textbf{Online FM}: 5-step Euler pseudo-$z_0$ introduces distribution mismatch; low LR can't correct biases, high LR causes patch-boundary artifacts.

\subsection{Data Scaling Law}

Fitting V6, V7, V9b to a log$_2$ model:
\begin{equation}
  \text{OOD CLIP \%} = 0.0115 \times \log_2(\text{prompts}) + 0.7744
  \label{eq:scaling}
\end{equation}
This predicted all three points with zero error (Figure~\ref{fig:scaling}).
However, V9c at 4{,}007 prompts yielded 88.8\% (predicted: 91.2\%), a $-$2.4pp miss.
The rank-64 architecture had saturated at $\sim$90\%.

\begin{figure}[h]
  \centering
  \includegraphics[width=0.7\textwidth]{figures/v9c_scaling_law.png}
  \caption{Data scaling law. The log$_2$ fit (dashed) predicts V6--V9b perfectly but fails at V9c (red X): 91.2\% predicted vs 88.8\% actual. The rank-64 ceiling is at $\sim$90\%.}
  \label{fig:scaling}
\end{figure}

\subsection{Breaking the Ceiling: Rank-128}

V10 doubled LoRA rank to 128 ($\sim$700M trainable parameters).
Fresh SVD initialization (rank-64 checkpoints incompatible) required 12{,}000 steps to converge, achieving 90.4\%---surpassing the rank-64 ceiling.

\begin{figure}[h]
  \centering
  \includegraphics[width=0.85\textwidth]{figures/v10_convergence.png}
  \caption{Per-prompt OOD CLIP for V9b (r64), V10 (r128, 6k cold), V10b (r128, 12k). V10's high variance (red) stabilizes with extended training (purple).}
  \label{fig:convergence}
\end{figure}

\section{Memory and Speed Analysis}

\subsection{Current Implementation (INT8 Storage)}

Our TernaryLinear stores ternary weights as INT8 (1 byte/weight) and casts to BF16 during forward via \texttt{F.linear(x, W\_q.to(x.dtype))}.
Table~\ref{tab:memory} reports measured inference metrics.

\begin{table}[h]
\centering
\caption{Inference benchmarks at 1024$\times$1024, 30 steps on A100-80GB. Speed is wall-clock time per image.}
\label{tab:memory}
\small
\begin{tabular}{lrrrrr}
\toprule
\textbf{Config} & \textbf{Static} & \textbf{Peak} & \textbf{Peak \%} & \textbf{Speed} & \textbf{Rel.} \\
\midrule
BF16 (baseline) & 31.44\,GB & 33.83\,GB & 100\% & 15.1\,s & 1.00$\times$ \\
Ternary r64 (INT8) & 21.02\,GB & 23.41\,GB & 69.2\% & 19.7\,s & 0.77$\times$ \\
Ternary r128 (INT8) & 25.10\,GB & 27.48\,GB & 81.2\% & 19.9\,s & 0.76$\times$ \\
\midrule
Ternary r64 (packed) & 24.45\,GB & 26.83\,GB & 79.3\% & 34.8\,s & 0.43$\times$ \\
Ternary r128 (packed) & 25.74\,GB & 28.13\,GB & 83.1\% & 35.0\,s & 0.43$\times$ \\
\bottomrule
\end{tabular}
\end{table}

\textbf{INT8 storage} yields a 30.8\% peak VRAM reduction (rank-64) at 23\% speed overhead.
The overhead comes from the INT8$\to$BF16 cast and extra LoRA matmuls.

\textbf{Rank-128} costs +4\,GB peak vs rank-64 (81.2\% vs 69.2\% of BF16) for +0.4pp quality.

\subsection{2-Bit Packed Storage (Attempted)}

We implemented 2-bit packing (4 ternary weights per byte) to achieve the theoretical 8$\times$ compression.
The packed representation reduces transformer weight buffers from 11.08\,GB to 2.77\,GB ($-$75\%).

\textbf{However}, the Python-level unpacking (\texttt{torch.stack}, reshape, subtract) creates temporary tensors that negate the static savings, resulting in \emph{higher} peak VRAM (79.3\% vs 69.2\%) and 2.3$\times$ slower inference.

\subsection{Custom Triton Kernel (Attempted)}

We wrote a Triton kernel performing on-the-fly 2-bit unpacking during matrix multiplication.
The kernel was \textbf{30--50$\times$ slower} than cuBLAS because:
\begin{itemize}[nosep,leftmargin=1.5em]
  \item Ternary unpacking forces column-by-column outer products instead of tiled \texttt{tl.dot}
  \item cuBLAS uses Tensor Cores optimized for BF16 GEMM; our kernel uses scalar FP32
  \item A100's memory bandwidth (2\,TB/s) already makes INT8$\to$BF16 cast nearly free
\end{itemize}

We also evaluated \textbf{BitBLAS} (Microsoft's INT2 library) but it lacked precompiled A100 kernels.
\textbf{GemLite} and \textbf{TorchAO} support 2-bit but with FP16 (not BF16) activations.

\subsection{The Kernel Gap}

The honest assessment: without \textbf{fused ternary matmul kernels} that directly operate on 2-bit packed weights with BF16 activations, the memory benefit is limited to INT8 storage ($\sim$30\% reduction).
The paper's claimed 5.1$\times$ memory reduction and latency improvements require custom CUDA kernels that have not been released and do not exist in any publicly available library for BF16 activations on A100.

\section{Quality Results}

\begin{table}[h]
\centering
\caption{Best results on 20 OOD diverse prompts.}
\label{tab:quality}
\small
\begin{tabular}{lccccc}
\toprule
\textbf{Model} & \textbf{Rank} & \textbf{CLIP} & \textbf{Aes.} & \textbf{LPIPS$\downarrow$} & \textbf{vs BF16} \\
\midrule
BF16 & --- & 0.3375 & 5.842 & ref & 100\% \\
V9b & 64 & 0.3036 & \textbf{5.939} & \textbf{0.664} & 90.0\% \\
V10b & 128 & \textbf{0.3050} & 5.686 & 0.719 & \textbf{90.4\%} \\
\bottomrule
\end{tabular}
\end{table}

V10b achieves the highest CLIP alignment while V9b retains better aesthetic quality.
The rank-128 cold-start path produces semantically stronger but perceptually rougher images.

\begin{figure}[h]
  \centering
  \includegraphics[width=0.85\textwidth]{figures/v10_highlights.png}
  \caption{Visual comparison: BF16 $|$ V9b (r64) $|$ V10b (r128) for 8 selected prompts. Color bars show CLIP \% of BF16. V10b's p16 (magic forest) exceeds BF16 at 103.7\%.}
  \label{fig:highlights}
\end{figure}

\section{Discussion}

\paragraph{Data vs capacity.}
The scaling law breakdown at rank-64 reveals a fundamental tension: low-rank adapters have finite capacity to represent diverse velocity fields.
With $\sim$350M trainable parameters, the model saturates at $\sim$2{,}000 prompts.
Rank-128 ($\sim$700M params) pushes this ceiling higher, but at 2$\times$ parameter cost.

\paragraph{Cold-start penalty.}
Rank-128 requires 2$\times$ more training steps from SVD init (12k vs 6k) and still shows under-converged prompts (p15 astronaut: 58.5\%).
A partial warm-start strategy---copying scale parameters from rank-64 and SVD-initializing only the extra LoRA dimensions---could accelerate convergence.

\paragraph{The kernel bottleneck.}
Current GPU software stacks (cuBLAS, Triton, PyTorch) are not designed for ternary arithmetic.
The 1.58-bit FLUX paper's efficiency claims rely on unreleased custom kernels.
Until fused ternary matmul kernels supporting BF16 activations become available (via BitBLAS updates, TorchAO maturation, or custom CUDA), the practical deployment benefit is limited to INT8 storage savings.

\section{Conclusion}

We reproduce ternary FLUX, achieving 90.4\% of BF16 quality with 30.8\% peak VRAM reduction (rank-64) or 18.8\% (rank-128).
The quality recovery demonstrates that flow-matching distillation with offline teacher latents is effective for extreme quantization.
However, realizing the full memory and speed benefits of 1.58-bit quantization requires custom kernel engineering that remains an open challenge.

\paragraph{Artifacts.}
Code: \url{https://github.com/ugonfor/1.58bit-flux}.
Weights and data: \url{https://huggingface.co/ugonfor/1.58bit-flux-reproduction}.

\begin{thebibliography}{9}
\bibitem{yang2024} C.~Yang et al., ``1.58-bit FLUX,'' arXiv:2412.18653, 2024.
\bibitem{flux} Black Forest Labs, ``FLUX.1,'' \url{https://github.com/black-forest-labs/flux}, 2024.
\bibitem{lipman2023} Y.~Lipman et al., ``Flow Matching for Generative Modeling,'' ICLR, 2023.
\bibitem{hu2022} E.~J.~Hu et al., ``LoRA: Low-Rank Adaptation of Large Language Models,'' ICLR, 2022.
\bibitem{zhang2018} R.~Zhang et al., ``The Unreasonable Effectiveness of Deep Features as a Perceptual Metric,'' CVPR, 2018.
\bibitem{radford2021} A.~Radford et al., ``Learning Transferable Visual Models From Natural Language Supervision,'' ICML, 2021.
\bibitem{bitblas} Microsoft, ``BitBLAS,'' \url{https://github.com/microsoft/BitBLAS}, 2024.
\end{thebibliography}

\newpage
\appendix
\section{Full 20-Prompt Grid}

\begin{figure}[H]
  \centering
  \includegraphics[width=0.6\textwidth]{figures/v10_full_grid.png}
  \caption{All 20 prompts: BF16 $|$ V9b $|$ V10b with per-prompt CLIP \% bars.}
\end{figure}

\end{document}
